\documentclass[aspectratio=169, 10pt]{beamer}

% --- Typography & Vietnamese Support ---
\usepackage{fontspec}
\setmainfont{Arial}
\setsansfont{Arial}
\setmonofont{Consolas}[Scale=0.88]

\usefonttheme{professionalfonts}

% --- Color Palette (Matching the Handbook) ---
\usepackage{xcolor}
\definecolor{primaryblue}{RGB}{10, 47, 100}      % NEU Deep Navy
\definecolor{secondarygreen}{RGB}{0, 128, 128}   % Teal / Modern Tech Green
\definecolor{accentorange}{RGB}{217, 119, 6}     % Warm Amber
\definecolor{darkgray}{RGB}{45, 55, 72}
\definecolor{lightbg}{RGB}{247, 250, 252}
\definecolor{boxbg}{RGB}{240, 244, 248}
\definecolor{codebg}{RGB}{245, 245, 245}

% --- Theme Setup ---
\setbeamercolor{normal text}{fg=darkgray, bg=white}
\setbeamercolor{structure}{fg=primaryblue}
\setbeamercolor{frametitle}{fg=white, bg=primaryblue}
\setbeamercolor{title}{fg=white}
\setbeamercolor{subtitle}{fg=accentorange}
\setbeamercolor{item}{fg=secondarygreen}
\setbeamercolor{subitem}{fg=primaryblue}

% --- Tcolorbox for Custom Slide Containers ---
\usepackage{tcolorbox}
\tcbuselibrary{skins}

\newtcolorbox{cbox}[2][primaryblue]{
  colback=#1!4,
  colframe=#1,
  coltitle=white,
  fonttitle=\bfseries\footnotesize,
  title={#2},
  arc=3pt,
  boxrule=0.9pt,
  left=4pt, right=4pt, top=3pt, bottom=3pt, boxsep=1pt,
  toptitle=1.5pt, bottomtitle=1.5pt,
  enhanced,
  drop fuzzy shadow=black!10
}

\newtcolorbox{promptbox}[1]{
  colback=lightbg,
  colframe=primaryblue!60,
  fonttitle=\bfseries\footnotesize,
  coltitle=primaryblue!90,
  title={#1},
  arc=3pt,
  boxrule=0.7pt,
  left=5pt, right=5pt, top=3pt, bottom=3pt, boxsep=1pt,
  enhanced
}

% --- TikZ Setup ---
\usepackage{tikz}
\usetikzlibrary{shapes.geometric, arrows.meta, positioning, calc, shadows}

% --- Tables & Code ---
\usepackage{booktabs}
\usepackage{tabularx}
\usepackage{listings}
\usepackage{xurl}

\lstset{
  backgroundcolor=\color{codebg},
  basicstyle=\ttfamily\scriptsize,
  breaklines=true,
  frame=single,
  rulecolor=\color{gray!30},
  keywordstyle=\color{primaryblue}\bfseries,
  commentstyle=\color{gray},
  stringstyle=\color{accentorange}
}

% --- Beamer Template Customization ---
\setbeamertemplate{navigation symbols}{}

\setbeamertemplate{frametitle}{%
  \nointerlineskip%
  \begin{beamercolorbox}[wd=\paperwidth,ht=3.2ex,dp=1.2ex,leftskip=0.5cm,rightskip=0.5cm]{frametitle}%
    \textbf{\insertframetitle}%
  \end{beamercolorbox}%
}

\setbeamertemplate{footline}{%
  \begin{beamercolorbox}[wd=\paperwidth,ht=2.4ex,dp=1.0ex,leftskip=0.5cm,rightskip=0.5cm]{author in head/foot}%
    \tiny\color{gray}%
    TS. Vũ Đức Minh (NEU) \hfill DSAI1003: Các nguyên lý lập trình với Python \hfill \insertframenumber{} / \inserttotalframenumber
  \end{beamercolorbox}%
}

% --- Metadata ---
\title{CẨM NANG HỌC PHẦN \& PHƯƠNG PHÁP TỰ HỌC}
\subtitle{DSAI1003: CÁC NGUYÊN LÝ LẬP TRÌNH CƠ BẢN VỚI PYTHON}
\author{TS. Vũ Đức Minh}
\institute{Khoa Khoa học Dữ liệu và Trí tuệ Nhân tạo -- Trường Công nghệ \\ Trường Đại học Kinh tế Quốc dân (NEU)}
\date{Học kỳ I, Năm học 2024 -- 2025}

\begin{document}

% ==============================================================================
% SLIDE 1: TITLE SLIDE
% ==============================================================================
\begin{frame}[plain]
  \vspace{0.2cm}
  \begin{tcolorbox}[colback=primaryblue, colframe=primaryblue, arc=4pt, left=12pt, right=12pt, top=10pt, bottom=10pt]
    \centering
    \textcolor{white}{\footnotesize\textbf{TRƯỜNG ĐẠI HỌC KINH TẾ QUỐC DÂN -- TRƯỜNG CÔNG NGHỆ}}\\[1pt]
    \textcolor{white}{\scriptsize\textbf{KHOA KHOA HỌC DỮ LIỆU VÀ TRÍ TUỆ NHÂN TẠO}}\\[8pt]
    {\large\textbf{\textcolor{white}{CẨM NANG HỌC PHẦN \& PHƯƠNG PHÁP TỰ HỌC}}}\\[3pt]
    {\large\textbf{\textcolor{accentorange}{DSAI1003: CÁC NGUYÊN LÝ LẬP TRÌNH CƠ BẢN VỚI PYTHON}}}\\[4pt]
    {\footnotesize\textcolor{white!80}{\textit{Fundamental Programming Concepts in Python}}}
  \end{tcolorbox}
  \vspace{0.15cm}
  \begin{tcolorbox}[colback=white, colframe=primaryblue!40, arc=3pt, left=8pt, right=8pt, top=5pt, bottom=5pt]
    \footnotesize
    \begin{tabularx}{\textwidth}{lX lX}
      \textbf{Giảng viên:} & TS. Vũ Đức Minh & \textbf{Chương trình:} & Khoa học dữ liệu (DSFE) \\
      \textbf{Khối lượng:} & 3 Tín chỉ (45h LT + 22.5h TH) & \textbf{Quyết định:} & 975/QĐ-ĐHKTQD (30/03/2024) \\
       \textbf{Website:} & \url{https://vudmvn.github.io/Fundamental-of-Programming-Concepts-in-Python/} \\
    \end{tabularx}
  \end{tcolorbox}
\end{frame}

% ==============================================================================
% SLIDE 2: AGENDA
% ==============================================================================
\begin{frame}{Nội dung Báo cáo / Cẩm nang Môn học}
  \begin{columns}[T]
    \begin{column}{0.49\textwidth}
      \begin{cbox}[primaryblue]{Phần 1: Tổng quan Học phần (Syllabus)}
        \footnotesize
        \begin{itemize}\setlength{\itemsep}{1pt}
          \item Thông tin chung \& Triết lý đào tạo
          \item Mục tiêu môn học (COs G1 -- G7)
          \item Giáo trình chuẩn quốc tế \& Đánh giá
          \item Lộ trình đào tạo chi tiết 15 tuần
        \end{itemize}
      \end{cbox}
      \vspace{0.1cm}
      \begin{cbox}[secondarygreen]{Phần 2: Phương pháp Tự học Hiệu quả}
        \footnotesize
        \begin{itemize}\setlength{\itemsep}{1pt}
          \item 4 sai lầm kinh điển \& Nguyên tắc cốt lõi
          \item Chu trình tự học chủ động 9 bước
          \item Predict $\rightarrow$ Run $\rightarrow$ Explain \& Teach-Back
          \item Debugging 8 bước \& Bảng Error Log
        \end{itemize}
      \end{cbox}
    \end{column}
    \begin{column}{0.49\textwidth}
      \begin{cbox}[accentorange]{Phần 3: Sử dụng AI Hỗ trợ Đột phá}
        \footnotesize
        \begin{itemize}\setlength{\itemsep}{1pt}
          \item Định vị: AI Personal Tutor \& Coach
          \item Khung Prompt chuẩn mực 4 thành phần
          \item Nguyên tắc 5T \& Hệ gợi ý 3 cấp độ
          \item Bộ công cụ Prompt Toolkit \& Thẩm định
        \end{itemize}
      \end{cbox}
      \vspace{0.1cm}
      \begin{cbox}[darkgray]{Phần 4: Kế hoạch Hành động Tuần 1}
        \footnotesize
        \begin{itemize}\setlength{\itemsep}{1pt}
          \item Cài đặt môi trường (VS Code, Colab)
          \item Khám phá tài nguyên trên GitHub
          \item Khởi tạo Nhật ký học tập (Learning Log)
          \item Cam kết văn hóa học tập chủ động
        \end{itemize}
      \end{cbox}
    \end{column}
  \end{columns}
\end{frame}

% ==============================================================================
% PART 1: SYLLABUS
% ==============================================================================
\begin{frame}[plain]
  \vfill
  \begin{tcolorbox}[colback=primaryblue, colframe=primaryblue, arc=5pt, left=15pt, right=15pt, top=15pt, bottom=15pt]
    \textcolor{accentorange}{\large\textbf{PHẦN 1}}\\[4pt]
    {\LARGE\textbf{\textcolor{white}{TỔNG QUAN HỌC PHẦN (SYLLABUS)}}}\\[6pt]
    \textcolor{white!80}{\small Thông tin môn học, Mục tiêu, Giáo trình, Đánh giá \& Lộ trình 15 tuần}
  \end{tcolorbox}
  \vfill
\end{frame}

% SLIDE 4: Course Introduction & Philosophy
\begin{frame}{1. Giới thiệu Học phần \& Triết lý Tiếp cận}
  \begin{columns}[T]
    \begin{column}{0.49\textwidth}
      \begin{cbox}[primaryblue]{Vị trí môn học trong CTĐT}
        \small
        \begin{itemize}\setlength{\itemsep}{2pt}
          \item \textbf{Mã học phần:} DSAI1003 (3 tín chỉ).
          \item \textbf{Đối tượng:} Sinh viên Cử nhân \textit{Khoa học dữ liệu trong Tài chính \& TMĐT (DSFE)}.
          \item \textbf{Tính chất:} Học phần \textbf{Bắt buộc cơ sở ngành} tiên quyết, đặt nền tảng tư duy thuật toán.
        \end{itemize}
      \end{cbox}
    \end{column}
    \begin{column}{0.49\textwidth}
      \begin{cbox}[secondarygreen]{Triết lý tiếp cận môn học}
        \small
        \begin{itemize}\setlength{\itemsep}{2pt}
          \item \textbf{Không đòi hỏi kinh nghiệm trước:} Chỉ cần nắm chắc đại số và tư duy logic.
          \item \textbf{Gắn liền ứng dụng thực tế:} Case study kinh tế, tài chính, TMĐT.
          \item \textbf{Từ bản chất đến nâng cao:} Thiết kế giải thuật $\rightarrow$ Cú pháp $\rightarrow$ OOP.
        \end{itemize}
      \end{cbox}
    \end{column}
  \end{columns}
  \vspace{0.12cm}
  \begin{tcolorbox}[colback=primaryblue!8, colframe=primaryblue!60, arc=3pt, left=8pt, right=8pt, top=3pt, bottom=3pt]
    \footnotesize
    \textbf{Cam kết chuẩn đầu ra:} Đủ năng lực tự thiết kế và cài đặt chương trình hoàn chỉnh giải quyết các bài toán phân tích kinh doanh thực tế, sẵn sàng cho các môn học chuyên sâu về Máy học và Trí tuệ nhân tạo.
  \end{tcolorbox}
\end{frame}

% SLIDE 5: Course Objectives (COs)
\begin{frame}{2. Mục tiêu Học phần (Course Learning Outcomes)}
  \begin{columns}[T]
    \begin{column}{0.49\textwidth}
      \begin{cbox}[primaryblue]{Tư duy \& Kỹ năng Lập trình Cơ bản}
        \footnotesize
        \begin{itemize}\setlength{\itemsep}{3pt}
          \item \textbf{G1 (Tư duy thuật toán):} Biểu diễn giải thuật bằng mã giả (pseudocode) và lưu đồ (flowchart).
          \item \textbf{G2 (Nền tảng Python):} Làm chủ kiểu dữ liệu, biểu thức, cấu trúc rẽ nhánh, vòng lặp và hàm.
          \item \textbf{G3 (Cấu trúc dữ liệu):} Quản lý dữ liệu hiệu quả với List, Tuple, Set, Dictionary và chuỗi.
        \end{itemize}
      \end{cbox}
    \end{column}
    \begin{column}{0.49\textwidth}
      \begin{cbox}[secondarygreen]{Kỹ năng Thực chiến \& Ứng dụng DSFE}
        \footnotesize
        \begin{itemize}\setlength{\itemsep}{3pt}
          \item \textbf{G4 (Tệp \& Ngoại lệ):} Xử lý tệp (I/O) và ngoại lệ an toàn, tin cậy.
          \item \textbf{G5 (Lập trình Hướng đối tượng):} Mô hình hóa thực thể bài toán qua Class, Object, Tính kế thừa.
          \item \textbf{G6 (Công cụ hiện đại):} Sử dụng thành thạo VS Code, Git, Google Colab và AI hỗ trợ.
          \item \textbf{G7 (Ứng dụng thực tế):} Giải quyết bài toán phân tích dữ liệu kinh tế -- tài chính.
        \end{itemize}
      \end{cbox}
    \end{column}
  \end{columns}
\end{frame}

% SLIDE 6: Textbooks & Evaluation
\begin{frame}{3. Giáo trình Chuẩn Quốc tế \& Phương thức Đánh giá}
  \begin{columns}[T]
    \begin{column}{0.49\textwidth}
      \begin{cbox}[primaryblue]{3 Bộ Giáo trình Chuẩn Quốc tế}
        \footnotesize
        \begin{enumerate}\setlength{\itemsep}{3pt}
          \item \textbf{[1] Python for Everyone} (3rd Ed., 2019)\\
          \textit{Cay Horstmann, Rance Necaise} (Wiley)\\
          $\rightarrow$ \textbf{Giáo trình chính} về cú pháp \& giải thuật.
          \item \textbf{[2] Intro to Python for CS and Data Science} (2020)\\
          \textit{Paul Deitel, Harvey Deitel} (Pearson)\\
          $\rightarrow$ \textbf{Giáo trình chính} về dữ liệu và AI.
          \item \textbf{[3] The Python Workbook} (2nd Ed., 2019)\\
          \textit{Ben Stephenson} (Springer)\\
          $\rightarrow$ Sách bài tập tự luyện kèm lời giải.
        \end{enumerate}
      \end{cbox}
    \end{column}
    \begin{column}{0.49\textwidth}
      \begin{cbox}[secondarygreen]{Cơ cấu Đánh giá Kết quả (3 Đầu điểm)}
        \scriptsize
        \begin{tabularx}{\textwidth}{lX c}
          \toprule
          \textbf{Hình thức} & \textbf{Nội dung \& Công cụ} & \textbf{Tỷ lệ} \\
          \midrule
          \textbf{Bài tập về nhà} & Lý thuyết \& code (LMS, 2 tuần/bài) & \textbf{30\%} \\
          \textbf{Thi giữa kỳ} & Bài thi Tuần 7 (Bài 1 -- 6, giấy/máy) & \textbf{30\%} \\
          \textbf{Thi cuối kỳ} & Thực hành máy phòng Lab (Tuần 15) & \textbf{40\%} \\
          \bottomrule
        \end{tabularx}
      \end{cbox}
      \vspace{0.1cm}
      \begin{cbox}[accentorange]{Quy định Học phần}
        \footnotesize
        \begin{itemize}\setlength{\itemsep}{2pt}
          \item \textbf{Chuyên cần:} Tham dự tối thiểu \textbf{80\%} số buổi học.
          \item \textbf{Điều kiện đạt môn:} Điểm tổng kết $\ge$ \textbf{5.0 / 10.0}.
        \end{itemize}
      \end{cbox}
    \end{column}
  \end{columns}
\end{frame}

% SLIDE 7: 15-week Roadmap Phase 1
\begin{frame}{4. Kế hoạch Giảng dạy 15 Tuần -- Giai đoạn 1 (Tuần 1 -- 7)}
  \scriptsize
  \begin{tabularx}{\textwidth}{c l X l}
    \toprule
    \textbf{Tuần} & \textbf{Bài học \& Nội dung} & \textbf{Trọng tâm Kiến thức} & \textbf{Đánh giá} \\
    \midrule
    \textbf{1} & \textbf{Bài 1: Giới thiệu \& Thuật toán} & Máy tính, thông dịch Python, mã giả, lưu đồ & Setup môi trường \\
    \textbf{2} & \textbf{Bài 2: Kiểu dữ liệu \& Biểu thức} & Số, chuỗi, lệnh gán, biến, toán tử, module & \textbf{Bài tập \#1} \\
    \textbf{3} & \textbf{Bài 3: Cấu trúc ra quyết định} & Rẽ nhánh \texttt{if}, \texttt{if-else}, \texttt{if-elif-else}, Boolean & Nộp BT \#1 \\
    \textbf{4} & \textbf{Bài 4: Cấu trúc vòng lặp} & Vòng lặp \texttt{while}, \texttt{for}, lặp lồng nhau, chuỗi & \textbf{Bài tập \#2} \\
    \textbf{5} & \textbf{Bài 5: Thiết kế hàm (Function)} & Cú pháp \texttt{def}, tham số, return, phạm vi biến & Nộp BT \#2 \\
    \textbf{6} & \textbf{Bài 6: Cấu trúc List \& Tuple} & Mảng 1D, 2D, slicing, các hàm tích hợp & \textbf{Bài tập \#3} \\
    \textbf{7} & \textbf{Thi giữa kỳ (Midterm Exam)} & Ôn tập củng cố \& thi đánh giá năng lực & \textbf{Thi GK (30\%)} \\
    \bottomrule
  \end{tabularx}
\end{frame}

% SLIDE 8: 15-week Roadmap Phase 2
\begin{frame}{5. Kế hoạch Giảng dạy 15 Tuần -- Giai đoạn 2 (Tuần 8 -- 15)}
  \scriptsize
  \begin{tabularx}{\textwidth}{c l X l}
    \toprule
    \textbf{Tuần} & \textbf{Bài học \& Nội dung} & \textbf{Trọng tâm Kiến thức} & \textbf{Đánh giá} \\
    \midrule
    \textbf{8--9} & \textbf{Bài 7: Thao tác Tệp \& Ngoại lệ} & Đọc/ghi file text, \texttt{sys.argv}, \texttt{try-except} & \textbf{Bài tập \#4} \\
    \textbf{10} & \textbf{Bài 8: Tập hợp (Set) \& Từ điển} & Key-value, bảng băm, cấu trúc JSON lồng nhau & Case study TMĐT \\
    \textbf{11--12} & \textbf{Bài 9: Đối tượng \& Lớp (OOP 1)} & Tư duy hướng đối tượng, \texttt{class}, \texttt{\_\_init\_\_}, method & Mini-lab OOP \\
    \textbf{13--14} & \textbf{Bài 10: Kế thừa \& Đa hình (OOP 2)} & Cây kế thừa, \texttt{super()}, ghi đè method, đa hình & Đề cương CK \\
    \textbf{15} & \textbf{Tổng kết \& Hướng dẫn thi} & Hệ thống hóa toàn bộ kiến thức môn học & Mock test phòng máy \\
    \bottomrule
  \end{tabularx}
\end{frame}

% ==============================================================================
% PART 2: EFFECTIVE SELF-STUDY METHODOLOGY
% ==============================================================================
\begin{frame}[plain]
  \vfill
  \begin{tcolorbox}[colback=secondarygreen, colframe=secondarygreen, arc=5pt, left=15pt, right=15pt, top=15pt, bottom=15pt]
    \textcolor{white}{\large\textbf{PHẦN 2}}\\[4pt]
    {\LARGE\textbf{\textcolor{white}{PHƯƠNG PHÁP TỰ HỌC LẬP TRÌNH HIỆU QUẢ}}}\\[6pt]
    \textcolor{white!80}{\small Tránh bẫy ảo tưởng năng lực, Chu trình 9 bước, Học tập chủ động \& Kỹ năng Debugging}
  \end{tcolorbox}
  \vfill
\end{frame}

% SLIDE 10: Core Principle & 4 Pitfalls
\begin{frame}{6. Bản chất Học Lập trình \& 4 Ngộ nhận Tai hại}
  \begin{tcolorbox}[colback=accentorange!12, colframe=accentorange, arc=3pt, left=6pt, right=6pt, top=2.5pt, bottom=2.5pt]
    \centering\footnotesize\bfseries
    "Học lập trình bằng cách đọc ít hơn, thực hành nhiều hơn và tự suy nghĩ trước khi xem lời giải."
  \end{tcolorbox}
  \vspace{0.06cm}
  \begin{columns}[T]
    \begin{column}{0.49\textwidth}
      \begin{cbox}[accentorange]{Cái bẫy 1: Đọc lướt không chạy code}
        \scriptsize
        Chỉ nhìn mắt thấy code mẫu và tưởng mình đã hiểu. Khi tắt máy đi, hoàn toàn không gõ lại được một dòng.
      \end{cbox}
      \vspace{0.06cm}
      \begin{cbox}[accentorange]{Cái bẫy 2: Copy-paste mù quáng}
        \scriptsize
        Copy code trên mạng hoặc AI, code chạy được nhưng không thể giải thích tại sao câu lệnh đó tồn tại.
      \end{cbox}
    \end{column}
    \begin{column}{0.49\textwidth}
      \begin{cbox}[accentorange]{Cái bẫy 3: Xem đáp án quá sớm}
        \scriptsize
        Vừa gặp bài khó là mở giải ngay. Bỏ qua bước "vật lộn" rèn luyện tư duy giải quyết vấn đề của não bộ.
      \end{cbox}
      \vspace{0.06cm}
      \begin{cbox}[accentorange]{Cái bẫy 4: Học dồn dập nhiều chủ đề}
        \scriptsize
        Cố học từ biến đến OOP trong 2 ngày mà không có thời gian làm bài tập để kiến thức lắng đọng.
      \end{cbox}
    \end{column}
  \end{columns}
\end{frame}

% SLIDE 11: 9-Step Active Study Cycle
\begin{frame}{7. Chu trình Tự học Lập trình Chủ động 9 Bước}
  \begin{center}
  \begin{tikzpicture}[
      stepb/.style={rectangle, rounded corners=3pt, draw=primaryblue, fill=primaryblue!8, text=primaryblue, font=\tiny\bfseries, minimum height=0.58cm, minimum width=2.3cm, align=center},
      stepg/.style={rectangle, rounded corners=3pt, draw=secondarygreen, fill=secondarygreen!8, text=secondarygreen, font=\tiny\bfseries, minimum height=0.58cm, minimum width=2.3cm, align=center},
      arr/.style={-{Stealth[scale=0.7]}, thick, primaryblue!80}
  ]
    % Row 1
    \node[stepb] (s1) at (0.6, 1.8) {1. Đọc (Read)};
    \node[stepb] (s2) at (4.0, 1.8) {2. Giải thích (Explain)};
    \node[stepb] (s3) at (7.4, 1.8) {3. Dự đoán (Predict)};
    \node[stepg] (s4) at (10.8, 1.8) {4. Chạy code (Run)};

    % Row 2
    \node[stepg] (s5) at (10.8, 0.65) {5. Tùy biến (Modify)};
    \node[stepg] (s6) at (7.4, 0.65) {6. Tự làm BT (Practice)};
    \node[stepg] (s7) at (4.0, 0.65) {7. Sửa lỗi (Debug)};
    \node[stepb] (s8) at (0.6, 0.65) {8. Tự đánh giá (Review)};

    % Step 9
    \node[stepb] (s9) at (0.6, -0.45) {9. Đúc kết (Reflect / Log)};

    % Connections
    \draw[arr] (s1) -- (s2); \draw[arr] (s2) -- (s3); \draw[arr] (s3) -- (s4);
    \draw[arr] (s4) -- (s5);
    \draw[arr] (s5) -- (s6); \draw[arr] (s6) -- (s7); \draw[arr] (s7) -- (s8);
    \draw[arr] (s8) -- (s9);
    \draw[arr, dashed] (s9.west) -- ++(-0.4, 0) |- (s1.west);
  \end{tikzpicture}
  \end{center}
  \vspace{0.05cm}
  \begin{tcolorbox}[colback=secondarygreen!8, colframe=secondarygreen, arc=3pt, left=8pt, right=8pt, top=3pt, bottom=3pt]
    \textbf{\footnotesize Phân bổ thời gian tự học vàng (Quy tắc 70/30):} \\
    \footnotesize Dành tối đa \textbf{30\% thời gian} cho việc đọc lý thuyết; hơn \textbf{70\% thời gian} phải dành trực tiếp cho: Gõ code, quan sát lỗi, phân tích traceback và giải bài tập độc lập!
  \end{tcolorbox}
\end{frame}

% SLIDE 12: Active Learning Techniques
\begin{frame}{8. Kỹ thuật Học tập Chủ động Đột phá}
  \begin{columns}[T]
    \begin{column}{0.49\textwidth}
      \begin{cbox}[primaryblue]{Predict $\rightarrow$ Run $\rightarrow$ Explain}
        \footnotesize
        \textbf{Trước khi bấm nút Run trên IDE:}
        \begin{itemize}\setlength{\itemsep}{1.5pt}
          \item Viết ra nháp chính xác màn hình sẽ in ra gì.
          \item Dự đoán giá trị biến trong bộ nhớ thay đổi ra sao.
          \item Nếu kết quả khác dự đoán $\rightarrow$ Đây là lúc học tốt nhất!
        \end{itemize}
      \end{cbox}
      \vspace{0.06cm}
      \begin{cbox}[secondarygreen]{Teach-Back (Giảng ngược lại)}
        \footnotesize
        Thử giải thích khái niệm cho người chưa biết gì. Nếu bạn chỉ nói \textit{"Nó chạy thế vì nó phải thế"} $\rightarrow$ Bạn chưa hiểu bản chất!
      \end{cbox}
    \end{column}
    \begin{column}{0.49\textwidth}
      \begin{cbox}[accentorange]{Active Recall (Gợi nhớ chủ động)}
        \footnotesize
        \textbf{Đừng đọc lại slide nhiều lần!}
        \begin{itemize}\setlength{\itemsep}{1.5pt}
          \item Gấp tài liệu lại và tự trả lời:
          \begin{itemize}\scriptsize
            \item \texttt{input()} trả về kiểu gì?
            \item Ép kiểu số nguyên dùng hàm nào?
            \item Tại sao chuỗi + số lại lỗi?
          \end{itemize}
        \end{itemize}
      \end{cbox}
      \vspace{0.06cm}
      \begin{cbox}[darkgray]{Spaced Repetition (Lặp lại ngắt quãng)}
        \footnotesize
        Ôn tập theo chu kỳ giãn cách:\\
        \centering\textbf{Ngày 1 $\rightarrow$ Ngày 3 $\rightarrow$ Ngày 7 $\rightarrow$ Ngày 14}\\
        \raggedright Biến kiến thức tạm thời thành phản xạ tư duy dài hạn.
      \end{cbox}
    \end{column}
  \end{columns}
\end{frame}

% SLIDE 13: The 5-Level Learning Ladder
\begin{frame}{9. Nấc thang Luyện tập 5 Cấp độ (Learning Ladder)}
  \begin{center}
  \begin{tikzpicture}[
    ladder/.style={rectangle, rounded corners=3pt, text width=12.2cm, minimum height=0.60cm, inner sep=4pt, font=\scriptsize, align=left}
  ]
    \node[ladder, fill=accentorange!30, draw=accentorange] at (0, 2.6) {\textbf{Cấp 5 -- Mini Project:} Xây dựng ứng dụng console quản lý đơn hàng, phân tích tệp giao dịch.};
    \node[ladder, fill=secondarygreen!30, draw=secondarygreen] at (0, 1.95) {\textbf{Cấp 4 -- Bài toán thực tế:} Tính tiền điện lũy tiến, tính thuế TNCN, tính lãi suất kép ngân hàng.};
    \node[ladder, fill=secondarygreen!15, draw=secondarygreen] at (0, 1.3) {\textbf{Cấp 3 -- Kết hợp khái niệm:} Kết hợp đồng thời vòng lặp \texttt{for} và cấu trúc rẽ nhánh \texttt{if-else}.};
    \node[ladder, fill=primaryblue!20, draw=primaryblue] at (0, 0.65) {\textbf{Cấp 2 -- Viết lại từ mô tả:} Che lời giải mẫu, đọc đề bài và tự mình viết code từ con số 0.};
    \node[ladder, fill=primaryblue!10, draw=primaryblue] at (0, 0) {\textbf{Cấp 1 -- Sửa ví dụ có sẵn:} Đổi giá trị biến, đổi kiểu số nguyên sang số thực, quan sát kết quả.};
  \end{tikzpicture}
  \end{center}
  \vspace{0.1cm}
  \centering\small\itshape
  Không bao giờ nhảy cóc từ ví dụ cơ bản lên ngay bài toán phức tạp mà bỏ qua các nấc thang trung gian!
\end{frame}

% SLIDE 14: Debugging Framework
\begin{frame}{10. Kỹ năng Debugging \& Phân loại Lỗi Tiêu chuẩn}
  \begin{columns}[T]
    \begin{column}{0.49\textwidth}
      \begin{cbox}[primaryblue]{Quy trình 8 bước Debugging tiêu chuẩn}
        \scriptsize
        \begin{enumerate}\setlength{\itemsep}{1.5pt}
          \item Đọc toàn bộ lỗi từ dòng cuối ngược lên.
          \item Định vị chính xác số dòng gây ra lỗi.
          \item Nhận diện kiểu lỗi (\texttt{TypeError}, \texttt{IndexError}...).
          \item Dùng \texttt{print()} kiểm tra giá trị biến tức thời.
          \item Dùng hàm \texttt{type()} kiểm tra kiểu dữ liệu biến.
          \item So sánh output thực tế vs. output kỳ vọng.
          \item Thử nghiệm với bộ input tối giản nhất.
          \item \textbf{Chỉ sửa 1 yếu tố} duy nhất trong mỗi lần thử!
        \end{enumerate}
      \end{cbox}
    \end{column}
    \begin{column}{0.49\textwidth}
      \begin{cbox}[accentorange]{Phân biệt 3 nhóm lỗi lập trình}
        \footnotesize
        \begin{itemize}\setlength{\itemsep}{3pt}
          \item \textbf{Syntax Error (Lỗi cú pháp):} Viết sai ngữ pháp ngôn ngữ (quên dấu \texttt{:}, đóng ngoặc). Python phát hiện ngay khi nạp.
          \item \textbf{Runtime Error (Thời gian chạy):} Cú pháp đúng nhưng sập khi chạy (chia cho 0, ép kiểu chuỗi chữ thành số nguyên).
          \item \textbf{Logic Error (Lỗi ngữ nghĩa):} Code chạy êm ru không báo lỗi nhưng \textbf{kết quả tính sai}! (Nguy hiểm nhất).
        \end{itemize}
      \end{cbox}
    \end{column}
  \end{columns}
\end{frame}

% SLIDE 15: Error Log
\begin{frame}{11. Nhật ký Học tập \& Sổ tay Bắt lỗi (Error Log)}
  \begin{tcolorbox}[colback=secondarygreen!8, colframe=secondarygreen, arc=3pt, left=8pt, right=8pt, top=3pt, bottom=3pt]
    \textbf{\footnotesize Biến lỗi sai thành tài sản học tập cá nhân:} \\
    \footnotesize Mỗi lỗi sai gặp phải khi được ghi lại cẩn thận sẽ trở thành kinh nghiệm quý báu giúp bạn không bao giờ mắc lại trong phòng thi.
  \end{tcolorbox}
  \vspace{0.1cm}
  \scriptsize
  \begin{tabularx}{\textwidth}{l X X c}
    \toprule
    \textbf{Lỗi gặp phải} & \textbf{Nguyên nhân gốc rễ} & \textbf{Cách khắc phục} & \textbf{Hiểu?} \\
    \midrule
    \texttt{TypeError: can only} & \texttt{input()} luôn trả về chuỗi kí tự, không thể cộng & Ép kiểu: \texttt{int(input())} hoặc & $\checkmark$ \\
    \texttt{concatenate str to str} & trực tiếp số. & \texttt{float(input())}. & \\
    \midrule
    \texttt{IndexError: list index} & Chỉ số mảng từ $0$ đến $N-1$, duyệt đến $N$ gây & Sửa biên: \texttt{range(len(lst))}. & $\checkmark$ \\
    \texttt{out of range} & tràn chỉ số. & & \\
    \midrule
    Vòng lặp vô tận (\texttt{while}) & Quên lệnh tăng biến đếm ở cuối thân vòng lặp. & Bổ sung lệnh: \texttt{i += 1}. & $\checkmark$ \\
    \bottomrule
  \end{tabularx}
  \vspace{0.15cm}
  \centering\footnotesize\itshape
  Hãy duy trì tệp \texttt{learning\_log.md} sau mỗi buổi học để theo dõi sự tiến bộ hàng tuần!
\end{frame}

% ==============================================================================
% PART 3: LEARNING PYTHON WITH AI
% ==============================================================================
\begin{frame}[plain]
  \vfill
  \begin{tcolorbox}[colback=accentorange, colframe=accentorange, arc=5pt, left=15pt, right=15pt, top=15pt, bottom=15pt]
    \textcolor{primaryblue}{\large\textbf{PHẦN 3}}\\[4pt]
    {\LARGE\textbf{\textcolor{white}{CHIẾN LƯỢC SỬ DỤNG AI (CHATGPT) HỖ TRỢ HỌC TẬP}}}\\[6pt]
    \textcolor{white!90}{\small AI Tutor, 4-Part Prompt, 5T Principle, Prompt Toolkit \& Verification}
  \end{tcolorbox}
  \vfill
\end{frame}

% SLIDE 17: AI Positioning (Tutor vs. Crutch)
\begin{frame}{12. Định vị Vai trò Chuẩn xác của AI trong Học Lập trình}
  \begin{columns}[T]
    \begin{column}{0.49\textwidth}
      \begin{cbox}[accentorange]{CẢNH BÁO: Thói quen lạm dụng AI}
        \footnotesize
        \textbf{Quy trình sai lầm:}\\
        Đề bài $\rightarrow$ Giao AI $\rightarrow$ Chép code $\rightarrow$ Bế tắc!
        \vspace{2pt}
        \begin{itemize}\setlength{\itemsep}{2pt}
          \item Hoàn toàn không hiểu tại sao câu lệnh đó được dùng.
          \item Vào phòng thi không có AI $\rightarrow$ Hoàn toàn bế tắc.
          \item Kỹ năng tư duy logic bị thoái hóa theo thời gian.
        \end{itemize}
      \end{cbox}
    \end{column}
    \begin{column}{0.49\textwidth}
      \begin{cbox}[secondarygreen]{KHUYẾN NGHỊ: AI là Huấn luyện viên (Coach)}
        \footnotesize
        \textbf{Quy trình chuẩn mực:}\\
        Tự nghĩ $\rightarrow$ Thử code $\rightarrow$ Xin gợi ý $\rightarrow$ Tự làm!
        \vspace{2pt}
        \begin{itemize}\setlength{\itemsep}{2pt}
          \item Nhờ AI \textbf{giải thích} khái niệm khó bằng ngôn ngữ đời thường.
          \item Nhờ AI \textbf{gợi ý từng bước} (không xin trọn vẹn đáp án).
          \item Nhờ AI \textbf{đặt câu hỏi ngược lại} để kiểm tra độ hiểu.
          \item Nhờ AI \textbf{review code} do chính mình tự viết.
        \end{itemize}
      \end{cbox}
    \end{column}
  \end{columns}
\end{frame}

% SLIDE 18: 4-Part Prompt Framework
\begin{frame}{13. Khung Cấu trúc Prompt Chuẩn mực 4 Thành phần}
  \scriptsize
  \begin{tabularx}{\textwidth}{l p{3.2cm} X}
    \toprule
    \textbf{Thành phần} & \textbf{Ý nghĩa câu hỏi} & \textbf{Ví dụ minh họa thực tế} \\
    \midrule
    \textbf{1. Context} & Bạn là ai, cấp độ nào? & \textit{"Tôi là sinh viên năm nhất đang bắt đầu học lập trình Python cơ bản..."} \\
    \textbf{2. Problem} & Bạn đang vướng chỗ nào? & \textit{"Tôi chưa phân biệt được sự khác nhau giữa biến số và giá trị gán..."} \\
    \textbf{3. Task} & Muốn AI làm gì cụ thể? & \textit{"Hãy giải thích bằng 1 ví dụ đời thường và 1 đoạn code Python dưới 8 dòng..."} \\
    \textbf{4. Constraint} & Ràng buộc câu trả lời? & \textit{"Không dùng kiến thức hàm/OOP, đặt cho tôi 2 câu hỏi kiểm tra, chưa vội đưa đáp án."} \\
    \bottomrule
  \end{tabularx}
  \vspace{0.15cm}
  \begin{tcolorbox}[colback=secondarygreen!8, colframe=secondarygreen, arc=3pt, left=8pt, right=8pt, top=3pt, bottom=3pt]
    \textbf{\footnotesize Sự khác biệt:}\\
    \footnotesize So với câu hỏi cộc lốc \textit{"Biến trong Python là gì?"}, prompt có cấu trúc giúp AI cá nhân hóa câu trả lời đúng 100\% trình độ của bạn!
  \end{tcolorbox}
\end{frame}

% SLIDE 19: The 5T Principle
\begin{frame}{14. Nguyên tắc 5T khi Tương tác cùng AI}
  \begin{center}
  \begin{tikzpicture}[
    tbox/.style={rectangle, rounded corners=4pt, draw=primaryblue, fill=primaryblue!6, text=primaryblue, minimum height=0.9cm, minimum width=2.0cm, align=center, font=\scriptsize\bfseries},
    arr/.style={-{Stealth[scale=0.7]}, thick, secondarygreen}
  ]
    \node[tbox] (t1) at (0, 0) {1. Think\\[1pt]\tiny\mdseries Suy nghĩ độc lập};
    \node[tbox] (t2) at (2.6, 0) {2. Try\\[1pt]\tiny\mdseries Tự viết thử code};
    \node[tbox] (t3) at (5.2, 0) {3. Talk\\[1pt]\tiny\mdseries Hỏi đúng chỗ nghẽn};
    \node[tbox, draw=primaryblue!90, fill=primaryblue!15] (t4) at (7.8, 0) {4. Test\\[1pt]\tiny\mdseries Kiểm chứng trên máy};
    \node[tbox] (t5) at (10.4, 0) {5. Teach-back\\[1pt]\tiny\mdseries Tự đúc kết lại};

    \draw[arr] (t1) -- (t2);
    \draw[arr] (t2) -- (t3);
    \draw[arr] (t3) -- (t4);
    \draw[arr] (t4) -- (t5);
  \end{tikzpicture}
  \end{center}
  \vspace{0.05cm}
  \begin{columns}[T]
    \begin{column}{0.49\textwidth}
      \footnotesize
      \begin{itemize}\setlength{\itemsep}{2pt}
        \item \textbf{Think:} Xác định Input, Output, thuật toán trước khi mở khung chat.
        \item \textbf{Try:} Tự gõ bản phác thảo đầu tiên, dù có lỗi.
        \item \textbf{Talk:} Cung cấp đoạn code hiện có và hỏi đúng chỗ tắc.
      \end{itemize}
    \end{column}
    \begin{column}{0.49\textwidth}
      \footnotesize
      \begin{itemize}\setlength{\itemsep}{2pt}
        \item \textbf{Test:} Luôn chạy code trên IDE, thử các dữ liệu biên (số 0, số âm, mảng rỗng).
        \item \textbf{Teach-back:} Đóng khung chat lại và tự giải thích lại giải pháp bằng lời mình.
      \end{itemize}
    \end{column}
  \end{columns}
\end{frame}

% SLIDE 20: 3-Level Hint System
\begin{frame}{15. Hệ thống Gợi ý 3 Cấp độ (Three-Level Hint System)}
  \begin{tcolorbox}[colback=primaryblue!8, colframe=primaryblue, arc=3pt, left=6pt, right=6pt, top=2pt, bottom=2pt]
    \centering\footnotesize\bfseries Khi làm bài tập khó, hãy yêu cầu AI cấp gợi ý theo từng nấc -- Không xin đáp án trọn vẹn!
  \end{tcolorbox}
  \vspace{0.08cm}
  \begin{columns}[T]
    \begin{column}{0.32\textwidth}
      \begin{cbox}[primaryblue]{Cấp 1: Gợi ý Ý tưởng}
        \scriptsize
        \textbf{(Idea Hint)}\\[2pt]
        Chỉ phân tích đề, chỉ ra hướng tiếp cận logic toán học/thuật toán.\\
        \vspace{2pt}
        $\rightarrow$ \textbf{Tuyệt đối không có dòng code nào!}
      \end{cbox}
    \end{column}
    \begin{column}{0.32\textwidth}
      \begin{cbox}[secondarygreen]{Cấp 2: Gợi ý Cấu trúc}
        \scriptsize
        \textbf{(Structure Hint)}\\[2pt]
        Gợi ý câu lệnh nên dùng: vòng \texttt{for}, lồng \texttt{if}, toán tử chia lấy dư \texttt{\%}...\\
        \vspace{2pt}
        $\rightarrow$ \textbf{Định hình khung logic.}
      \end{cbox}
    \end{column}
    \begin{column}{0.32\textwidth}
      \begin{cbox}[accentorange]{Cấp 3: Khung sườn Code}
        \scriptsize
        \textbf{(Scaffold Hint)}\\[2pt]
        Đưa ra khung code có vị trí trống (\texttt{\# TODO}) để tự suy nghĩ điền nốt.\\
        \vspace{2pt}
        $\rightarrow$ \textbf{Giữ quyền làm chủ bài giải.}
      \end{cbox}
    \end{column}
  \end{columns}
  \vspace{0.1cm}
  \begin{tcolorbox}[colback=accentorange!10, colframe=accentorange, arc=3pt, left=6pt, right=6pt, top=2pt, bottom=2pt]
    \centering\scriptsize\bfseries
    Chiến lược làm bài: Hãy nỗ lực giải bài ở Gợi ý Cấp 1. Chỉ mở Cấp 2 và Cấp 3 khi thực sự bế tắc sau 15 phút suy nghĩ!
  \end{tcolorbox}
\end{frame}

% SLIDE 21: Prompt Toolkit 1
\begin{frame}{16. Prompt Toolkit Thực chiến (1): Khái niệm \& Gợi ý}
  \begin{columns}[T]
    \begin{column}{0.49\textwidth}
      \begin{cbox}[primaryblue]{Học khái niệm mới (Concept Prompt)}
        \scriptsize
        Tôi là người mới bắt đầu học Python. Hãy giải thích cho tôi khái niệm [Hàm đệ quy].\\[2pt]
        \textbf{Yêu cầu:}
        \begin{enumerate}\setlength{\itemsep}{1pt}
          \item Bản chất dễ hiểu kèm 1 ví dụ đời thực.
          \item Đoạn code mẫu ngắn (dưới 10 dòng).
          \item 1 lỗi phổ biến người mới hay mắc.
          \item Đặt 2 câu hỏi kiểm tra, chưa vội đưa đáp án.
        \end{enumerate}
      \end{cbox}
    \end{column}
    \begin{column}{0.49\textwidth}
      \begin{cbox}[secondarygreen]{Xin gợi ý bài tập (Socratic Hint)}
        \scriptsize
        Tôi đang tự làm bài: "[NỘI DUNG ĐỀ BÀI]"\\[2pt]
        Code tôi đã thử viết:\\
        \texttt{[DÁN ĐOẠN CODE CỦA BẠN]}\\[2pt]
        Tôi đang bị tắc ở bước [MÔ TẢ VƯỚNG MẮC].\\[2pt]
        \textbf{Yêu cầu:} TUYỆT ĐỐI KHÔNG viết code giải hộ. Hãy cho tôi GỢI Ý CẤP ĐỘ 1 về mặt ý tưởng để tôi tự suy nghĩ tiếp.
      \end{cbox}
    \end{column}
  \end{columns}
\end{frame}

% SLIDE 22: Prompt Toolkit 2
\begin{frame}{17. Prompt Toolkit Thực chiến (2): Debug \& Code Review}
  \begin{columns}[T]
    \begin{column}{0.49\textwidth}
      \begin{cbox}[accentorange]{Hỗ trợ Debugging (Tìm nguyên nhân)}
        \scriptsize
        Tôi viết code Python sau để: [MỤC TIÊU CỦA CODE]\\[2pt]
        Code hiện tại: \texttt{[DÁN CODE]}\\[2pt]
        Lỗi nhận được: \texttt{[DÁN THÔNG BÁO LỖI]}\\[2pt]
        \textbf{Yêu cầu:}
        \begin{enumerate}\setlength{\itemsep}{1pt}
          \item Giải thích nguyên nhân tại sao có lỗi này.
          \item Chỉ rõ dòng code đang tiềm ẩn vấn đề.
          \item Đặt cho tôi 1 câu hỏi gợi mở để tôi tự sửa, KHÔNG sửa lại toàn bộ code.
        \end{enumerate}
      \end{cbox}
    \end{column}
    \begin{column}{0.49\textwidth}
      \begin{cbox}[primaryblue]{Review \& Tối ưu mã nguồn (Code Review)}
        \scriptsize
        Tôi đã hoàn thành xong bài tập này và code chạy đúng: \texttt{[DÁN CODE ĐÃ CHẠY ĐÚNG]}\\[2pt]
        Hãy đóng vai trò Giảng viên Python:
        \begin{enumerate}\setlength{\itemsep}{1pt}
          \item Đánh giá đặt tên biến, tính rõ ràng, cấu trúc logic.
          \item Code có xử lý được các trường hợp biên (edge cases) không?
          \item Gợi ý 1 điểm có thể viết gọn hơn chuẩn Pythonic và giải thích lý do.
        \end{enumerate}
      \end{cbox}
    \end{column}
  \end{columns}
\end{frame}

% SLIDE 23: Prompt Toolkit 3
\begin{frame}{18. Prompt Toolkit Thực chiến (3): Tạo Quiz Ôn tập}
  \begin{cbox}[primaryblue]{Tạo Bộ đề Quiz tự kiểm tra kiến thức (Self-Quiz Prompt)}
    \footnotesize
    Tôi vừa học xong các chủ đề: \textbf{[List, Tuple, Slicing và Vòng lặp for]}.\\[3pt]
    Hãy tạo cho tôi một bài kiểm tra ngắn gồm 6 câu hỏi:
    \begin{itemize}\setlength{\itemsep}{2pt}
      \item \textbf{2 câu Trắc nghiệm:} 4 lựa chọn A, B, C, D (chỉ 1 đáp án đúng).
      \item \textbf{2 câu Dự đoán kết quả (Predict the Output):} Đoạn code 3--5 dòng không dùng hàm nâng cao.
      \item \textbf{2 câu Tìm lỗi sai (Find the Bug):} Đoạn code có đúng 1 lỗi cần phát hiện và sửa.
    \end{itemize}
    \vspace{2pt}
    \textbf{Ràng buộc:} KHÔNG hiển thị đáp án trước. Hãy chờ tôi gửi câu trả lời rồi mới chấm điểm và giải thích cặn kẽ từng câu!
  \end{cbox}
  \vspace{0.15cm}
  \begin{tcolorbox}[colback=secondarygreen!8, colframe=secondarygreen, arc=3pt, left=8pt, right=8pt, top=3pt, bottom=3pt]
    \textbf{\footnotesize Mẹo nâng cao: Quiz từ chính lỗi sai trong quá khứ}\\
    \footnotesize Sao chép các dòng lỗi trong \texttt{learning\_log.md} của bạn vào prompt và yêu cầu: \textit{"Hãy tạo 5 câu hỏi xoay quanh đúng những lỗi sai này để kiểm tra xem tôi đã thực sự khắc phục được chưa!"}
  \end{tcolorbox}
\end{frame}

% SLIDE 24: 7-Point AI Verification Checklist
\begin{frame}{19. Bảng kiểm 7 Tiêu chí Thẩm định Thông tin từ AI}
  \begin{columns}[T]
    \begin{column}{0.49\textwidth}
      \footnotesize
      \begin{itemize}\setlength{\itemsep}{3pt}
        \item[$\square$] \textbf{1. Tính tương thích:} Code AI đưa ra có dùng thư viện lạ hay kiến thức vượt phạm vi môn học không?
        \item[$\square$] \textbf{2. Chạy thử nghiệm:} Đã copy code vào IDE và chạy thử không có lỗi biên dịch/runtime chưa?
        \item[$\square$] \textbf{3. Khớp đề bài:} Kết quả in ra có đúng từng chữ so với yêu cầu đầu ra của bài toán không?
        \item[$\square$] \textbf{4. Kiểm tra biên:} Code có chạy đúng khi input là số 0, số âm, danh sách rỗng?
      \end{itemize}
    \end{column}
    \begin{column}{0.49\textwidth}
      \footnotesize
      \begin{itemize}\setlength{\itemsep}{3pt}
        \item[$\square$] \textbf{5. Thuật ngữ chuẩn:} Thuật ngữ AI giải thích có đồng nhất với giáo trình chính thức môn học không?
        \item[$\square$] \textbf{6. Tự giải thích từng dòng:} Bạn có thể chỉ tay vào từng dòng code và giải thích tại sao nó ở đó?
        \item[$\square$] \textbf{7. Tái tạo độc lập:} Nếu tắt AI đi, bạn có thể tự mình viết lại code này trong 10 phút không?
      \end{itemize}
    \end{column}
  \end{columns}
  \vspace{0.15cm}
  \begin{tcolorbox}[colback=accentorange!15, colframe=accentorange, arc=3pt, left=8pt, right=8pt, top=3pt, bottom=3pt]
    \centering\footnotesize
    \textbf{Quy tắc Vàng Bất di Bất dịch:}\\[2pt]
    \textbf{"TUYỆT ĐỐI KHÔNG BAO GIỜ NỘP HOẶC LƯU LẠI MÃ NGUỒN\\MÀ BẢN THÂN BẠN KHÔNG THỂ TỰ MÌNH GIẢI THÍCH TƯỜNG TẬN!"}
  \end{tcolorbox}
\end{frame}

% ==============================================================================
% PART 4: ACTION PLAN FOR WEEK 1 & CONCLUSION
% ==============================================================================
\begin{frame}{20. Kế hoạch Hành động Ngay trong Tuần Đầu tiên}
  \begin{columns}[T]
    \begin{column}{0.32\textwidth}
      \begin{cbox}[primaryblue]{1. Chuẩn bị Môi trường}
        \scriptsize
        \begin{itemize}\setlength{\itemsep}{2pt}
          \item Cài đặt \textbf{Python 3.11+} từ trang chủ \url{python.org}.
          \item Cài đặt trình soạn thảo \textbf{VS Code} kèm extension Python.
          \item Đăng nhập sẵn tài khoản \textbf{Google Colab} trên trình duyệt.
        \end{itemize}
      \end{cbox}
    \end{column}
    \begin{column}{0.32\textwidth}
      \begin{cbox}[secondarygreen]{2. Khám phá Học liệu}
        \scriptsize
        \begin{itemize}\setlength{\itemsep}{2pt}
          \item Bookmark kho GitHub học phần: \url{github.com/vudmvn/Python}.
          \item Đọc kỹ Đề cương chi tiết (\texttt{syllabus-vn.md}).
          \item Đọc trước Chương 1 giáo trình \textit{Python for Everyone}.
        \end{itemize}
      \end{cbox}
    \end{column}
    \begin{column}{0.32\textwidth}
      \begin{cbox}[accentorange]{3. Sổ tay Tự học}
        \scriptsize
        \begin{itemize}\setlength{\itemsep}{2pt}
          \item Tạo thư mục học tập \texttt{DSAI1003\_Python/} trên máy.
          \item Tạo tệp \texttt{learning\_log.md} ghi chép sau mỗi tuần.
          \item Lưu sẵn bộ khung Prompt Toolkit để tương tác AI văn minh.
        \end{itemize}
      \end{cbox}
    \end{column}
  \end{columns}
  \vspace{0.2cm}
  \begin{tcolorbox}[colback=primaryblue!5, colframe=primaryblue, arc=3pt, left=8pt, right=8pt, top=6pt, bottom=6pt]
    \centering\small
    \textbf{\textcolor{primaryblue}{CHÚC CÁC BẠN SINH VIÊN KHOA KH DỮ LIỆU \& TRÍ TUỆ NHÂN TẠO}}\\[3pt]
    \textbf{\textcolor{secondarygreen}{LÀM CHỦ TƯ DUY LẬP TRÌNH \& BỨC PHÁ NĂNG LỰC CÙNG KỶ NGUYÊN AI!}}
  \end{tcolorbox}
\end{frame}

% SLIDE 26: THANK YOU
\begin{frame}[plain]
  \vfill
  \centering
  {\Huge\textbf{\textcolor{primaryblue}{THANK YOU!}}}\\[8pt]
  {\Large\textbf{\textcolor{secondarygreen}{CẢM ƠN CÁC BẠN ĐÃ LẮNG NGHE}}}\\[15pt]

  \begin{tcolorbox}[colback=white, colframe=primaryblue, arc=4pt, width=0.82\textwidth, left=10pt, right=10pt, top=8pt, bottom=8pt]
    \centering\footnotesize
    \textbf{\normalsize TS. Vũ Đức Minh}\\[3pt]
    Khoa Khoa học Dữ liệu và Trí tuệ Nhân tạo -- Trường Công nghệ\\
     Đại học Kinh tế Quốc dân (NEU)\\[6pt]
   
    \textbf{Website:} \url{https://vudmvn.github.io/Fundamental-of-Programming-Concepts-in-Python/}
  \end{tcolorbox}
  \vfill
\end{frame}

\end{document}